\documentclass[
  onecolumn,
  aps,
  %prl,
  amsmath,
  amssymb,
  superscriptaddress,
  nofootinbib,
  floatfix
]{revtex4-2}
\usepackage{mathtools}%enhanced math capabilities
\usepackage{bm}
\usepackage{dsfont}
\usepackage{graphicx}
\usepackage{subfigure}
\usepackage{pifont}
\usepackage{textcomp}
\usepackage{gensymb}
\usepackage{accents}
\usepackage{upgreek}
\usepackage{array}
\usepackage{color}
\usepackage{verbatim}

\usepackage{hyperref}
\hypersetup{colorlinks=true, citecolor=blue, urlcolor=blue, linkcolor=blue}

\makeatletter
\newcommand*{\supplementarystart}{%
  \close@column@grid%
  \clearpage%
  \onecolumngrid%
  \setcounter{enumiv}{0} % resets counter for references
  \setcounter{equation}{0} % resets counter for equations
  \setcounter{figure}{0} % resets counter for figs
  \setcounter{table}{0} % resets counter for tables
  \setcounter{page}{1}
  \c@secnumdepth=4
  \renewcommand{\theequation}{s\arabic{equation}} % equations numbered with S...
  \renewcommand{\bibnumfmt}[1]{[s##1]} % bibtems [S...]
  \renewcommand{\@onlinecite}{s\citealp} % citations [S...]
  \renewcommand{\cite}[1]{{[}\onlinecite{##1}{]}}
  \renewcommand{\thefigure}{s\arabic{figure}}
  \renewcommand{\thetable}{s\Roman{table}}
  \renewcommand{\thepage}{s\arabic{page}}
}
\makeatother

\newcommand{\s}{\sum\limits}
\newcommand{\p}{\prod\limits}
\newcommand{\pa}{\partial}
\newcommand{\il}{\int\limits}
\newcommand{\be}{\begin{equation}}
\newcommand{\e}{\end{equation}}
\newcommand{\bes}{\begin{equation*}}
\newcommand{\es}{\end{equation*}}
\newcommand{\beml}{\begin{subequations}}
\newcommand{\eml}{\end{subequations}}
\newcommand{\beq}{\begin{eqnarray}}
\newcommand{\eq}{\end{eqnarray}}
\newcommand{\ba}{\begin{array}}
\newcommand{\ea}{\end{array}}
\newcommand{\bpm}{\begin{pmatrix}}
\newcommand{\epm}{\end{pmatrix}}
\newcommand{\bc}{\begin{cases}}
\newcommand{\ec}{\end{cases}}
\newcommand{\lt}{\left}
\newcommand{\rt}{\right}
\newcommand{\n}{\nonumber}
\newcommand{\la}{\langle}
\newcommand{\ra}{\rangle}
\newcommand{\ep}{\varepsilon}
\newcommand{\dd}{\displaystyle}
\newcommand{\bs}{\mathbf}
\newcommand{\bb}{\boldsymbol}
\newcommand{\h}{^\dagger}
\newcommand{\0}{^{\phantom{\dagger}}}
\newcommand{\one}{\openone}
\newcommand{\uvec}[1]{\hat{\bb{#1}}}
\newcommand{\para}{\parallel}
%\newcommand{\d}{\!d\mkern-1.5mu}

\renewcommand{\log}{\mathop{\mathrm{ln}}\nolimits}

\DeclareMathOperator{\var}{var}
\DeclareMathOperator{\tr}{Tr}
\DeclareMathOperator{\diag}{diag}
\DeclareMathOperator{\im}{Im}
\DeclareMathOperator{\re}{Re}
\DeclareMathOperator{\ctg}{cotan}
\DeclareMathOperator{\arcsh}{arcsinh}
\DeclareMathOperator{\rot}{rot}
\DeclareMathOperator{\sign}{sign}
\DeclareMathOperator{\dv}{div}

\begin{document}
	
\title{Spin echo in the presence of Hermitian pulse}	


\date{\today}
%\email[bmurzaliev@science.ru.nl]
\begin{abstract}

We considered spin echo in NV center applying Hermitian pulse
\end{abstract}
\keywords{spin echo, NV center}	
\maketitle
	
\section{Introduction}
A major challenge in quantum computing is a control and measurement of qubits. One of the most promising candidate for building block of quantum computing and network is NV center in a diamond.
\section{Model}


Our goal is to study spin echo with finite duration pulse. Some papers
showed that finite-duration pulse can lead to non-negligible effect.

We have a Hamiltonian in lab frame

\begin{equation}
{\cal H}=-\gamma_{e}H_{0}S_{z}+\gamma_{n}H_{0}I_{z}+A\boldsymbol{S}\cdot\boldsymbol{I}+x(t)\omega_{1}(t)(S_{x}\cos(\omega t)+S_{y}\sin(\omega t))+
\end{equation}

\[
+y(t)\omega_{1}(t)(-S_{x}\sin(\omega t)+S_{y}\cos(\omega t))
\]

First and second terms are Zeeman fields for electron and nuclear
species, third term is hyperfine coupling, fourth term is $X$-pulse,
fifth term is $Y$-pulse.

\begin{equation}
\omega_{1}(t)=\omega_{1}(1-0.956(t/T_{2})^{2})e^{-(t/T_{2})^{2}}
\end{equation}

$-1\leq x(t)\leq1,-1\leq y(t)\leq1$ - modulation of signals, $T_{2}$
is width of the signal

Nuclear Zeeman frequencies are much larger than the hyperfine interaction
($\gamma_{n}H_{0}\gg A$), so that the hyperfine interaction can be
separated into secular and nonsecular contributions,

\begin{equation}
A\boldsymbol{S}\cdot\boldsymbol{I}\approx AS_{z}(a_{\parallel}I_{z}+a_{\perp}I_{x})
\end{equation}

with $a_{\parallel}=A\cos\theta$ , $a_{\perp}=A\sin\theta$.

We have a Hamiltonian in rotating frame

\begin{equation}
{\cal H}={\cal H}_{0}+{\cal H}_{P}(t)
\end{equation}

\begin{equation}
{\cal H}_{0}=(\omega-\gamma_{e}H_{0})\hat{S}_{z}+a_{\perp}\hat{S}_{z}\hat{I}_{x}+a_{\parallel}S_{z}I_{z}+\gamma_{n}H_{0}I_{z}
\end{equation}

with $\Delta=(\omega-\gamma_{e}H_{0})$,

\begin{equation}
{\cal H}_{P}(t)=x(t)\omega_{1}(t)S_{x}+y(t)\omega_{1}(t)S_{y}
\end{equation}

We have evolution $(\tau-\pi_{x}-2\tau-\pi_{y}-2\tau-\pi_{-x}-2\tau-\pi_{-y}-\tau)^{N}$

\begin{equation}
U=U_{\pi/2}U_{block}^{N}U_{\pi/2}
\end{equation}

where

\begin{equation}
U_{block}=U({\cal H}_{0},\tau)U_{\pi_{x}}U({\cal H}_{0},2\tau)U_{\pi_{y}}U({\cal H}_{0},2\tau)U_{\pi_{-x}}U({\cal H}_{0},2\tau)U_{\pi_{-y}}U({\cal H}_{0},\tau)
\end{equation}

So, if we have ideal pulses $\omega_{1}(t)t_{w}=\pi$, we have these
Hamiltonians and evolution for different time intervals.

If we have non-ideal pulses with flip angle $\beta=\int_{-\infty}^{\infty}\omega_{1}(t)dt=0.925221\omega_{1}T_{2}=\pi-\epsilon$.

According {[}Andreas Brinkmann{]} , we have to divide our Hamiltonian
into two parts

\begin{equation}
{\cal H}={\cal H}_{A}+{\cal H}_{B}
\end{equation}

with conditions : 1) $U_{A}(t_{a},t_{b})=\exp(-i{\cal H}_{A}(t_{b}-t_{a}))$
can be estimated analytically and periodicity ${\cal H}_{A}(t+NT)={\cal H}_{A}(t)$,
2) $H_{B}(t)$ is small, $\max||{\cal H}_{B}T||\ll1$ , ${\cal H}_{B}(t+NT)={\cal H}_{B}(t)$.

In our case we have
\begin{equation}
{\cal H}_{A}=\Delta\hat{S}_{z}+a_{\perp}\hat{S}_{z}\hat{I}_{x}+a_{\parallel}S_{z}I_{z}+\gamma_{n}H_{0}I_{z}+\omega_{\pi}(x(t)S_{x}+y(t)S_{y})
\end{equation}

\begin{equation}
{\cal H}_{B}=\omega_{\varepsilon}(x(t)S_{x}+y(t)S_{y})
\end{equation}

Now we want to transform Hamiltonian ${\cal H}_{B}$ in toggling reference
frame $\tilde{{\cal H}}_{B}=U_{A}^{\dagger}{\cal H}_{B}U_{A}$
\end{document}

\section{Average Hamiltonian}






\end{document}